Many systems in engineering and in finance evolve over time and can only be observed through noisy measurements. A drone flies over a field while a GPS receiver reports its position with an error of several meters. A radar station tracks an aircraft but measures only a noisy range and bearing. A financial market moves from day to day, and the price at which an asset trades reflects an underlying trend only through a lot of random variation. In all these situations the quantity of interest, the state of the system, is hidden. What we actually hold is a stream of measurements, each of them related to the state but corrupted by noise, and the task is to recover the state from them as well as possible.

Two difficulties make this task interesting. The first is that the state does not stand still. By the time a measurement arrives, the system has already moved on, so a good method has to combine what the measurements say with what we know about how the system evolves. The second is that an estimate on its own is of little use without a statement of how reliable it is. A navigation system that places the drone at some position should also say whether the uncertainty is one meter or fifty, because decisions built on the estimate depend on that difference.

Bayesian filtering addresses both difficulties at once. The idea is to describe our knowledge of the hidden state not by a single number but by a probability distribution, and to update this distribution recursively over time. Each step has the same form. A prediction carries the current knowledge forward through a model of the dynamics, and an update corrects this prediction with the newest measurement. The result is a running estimate of the state together with its uncertainty, computed as the data arrives \cite[Ch.~1]{sarkka:bayesian_filtering}. When the dynamics and the measurements are linear and all noise is Gaussian, the recursion can be carried out exactly and in closed form. This special case is the Kalman filter. Outside the linear Gaussian case this exactness is lost, and the extended Kalman filter approximates the same recursion for non-linear models. One further ingredient is needed before the theory becomes usable. The models contain parameters, above all the noise covariances, that are rarely known in advance. They can be estimated from the data by maximum likelihood, and conveniently the filter itself produces the likelihood as a by-product while it runs.

The second theme of this thesis is what these methods deliver in practice, on data that does not come from a simulation. Financial return series are a natural test. Their average level and their typical size both change over time, which suggests hidden states worth filtering, and at the same time financial theory warns that prices are hard to predict, which promises an honest test rather than an easy success. The thesis therefore asks two questions. The first is how the hidden state of a system and the associated uncertainty can be estimated recursively from noisy observations, and under which assumptions this estimate is exact. The second is how well the resulting filters predict a real financial time series compared with simple baseline models. 

The thesis is organized as follows. Chapter~\ref{chap:basic} collects the probabilistic foundations, in particular the Gaussian distribution and the closure properties that make the later recursions work, and introduces state space models together with the drone example that accompanies the first chapters. Chapter~\ref{chap:bayesian} derives the Bayesian filtering equations, the two-step recursion at the heart of everything that follows, and explains the related smoothing problem. Chapter~\ref{chap:kalman} specializes the recursion to linear Gaussian models, which yields the Kalman filter, and studies in simulation how the process and measurement noise covariances shape its behavior. Chapter~\ref{chap:nonlinear} treats non-linear models, develops the extended Kalman filter, and mentions the unscented Kalman filter as an alternative. Chapter~\ref{chap:parameter} turns to parameter estimation and develops the energy function, the negative log-likelihood that the filter computes along the way, first for linear Gaussian models and then for non-linear ones. Chapter~\ref{chap:portfolio} brings everything together in an application to financial time series. Starting from a model that treats returns as independent draws, it builds a sequence of models of increasing sophistication, tracks their hidden states with the filters of the earlier chapters, and compares their one-step-ahead predictions on data the models have not seen.
